\subsection{Bounding field excursions along null geodesics with applications to cosmology --- arXiv:2510.01554}

\textbf{Authors:} Aidan Herderschee, Aron C. Wall

\paragraph{主な主張}
Raychaudhuri 方程式と Einstein 方程式を用い、他の物質場が null energy condition (NEC) を満たすという仮定の下で、null geodesic に沿うスカラー場の変位に対する field excursion bound (FEB)
\[
\left|\log\frac{\theta_2}{\theta_1}\right|
\geq
4\sqrt{\frac{2\pi}{D-2}}\,d(X_2,X_1)
\]
を導出した。ここで $d(X_2,X_1)$ は scalar-field space metric で測った geodesic distance であり、ポテンシャルの具体形には依存しない。FLRW 宇宙への応用では、field excursion が e-fold 数に対して線形に制限されることを示した \cite{Herderschee:2025FEB}。また NEC が破れる半古典的時空に対しては quantum expansion を用いた quantum FEB を提案するが、これは strengthened quantum focusing condition (SQFC) を仮定する conjectural extension である。

\paragraph{新規性}
大きなスカラー場変位を null congruence の gravitational focusing と直接結び付け、特定の inflationary potential や slow-roll dynamics に依存しないモデル非依存の上限を与えた点が新しい。Lyth bound が観測量から field excursion の下限を与えるのに対し、本結果は重力的 focusing から上限を与える。

\paragraph{位置付け}
Inflation や moduli dynamics における trans-Planckian field excursion を評価する一般的な幾何学的制約として位置付けられる。特に hyperbolic moduli space では座標変位ではなく geodesic distance が bound に現れるため、modular cosmology のように $ds^2\propto d\tau d\bar\tau/(\operatorname{Im}\tau)^2$ を持つ系では、大きな $\operatorname{Im}\tau$ の変化が対数的な canonical distance に対応する点が重要である。
